\section{证明}
\label{appendix-proof}
\begin{lemma}
  给定系统 $\mathfrak{S} = (\mathcal{X}, \mathcal{X}_0, f)$、轨迹 $\tau=\langle x_0, x_1, \dots \rangle \in \mathfrak{T}(\mathfrak{S})$、标记函数 $\mathfrak{L}:\mathcal{X} \rightarrow \Sigma$ 和 NBA $\mathcal{A}=(\Sigma, Q, q_0, \delta, Acc)$。若能够找到相对于 $q_k \in Acc^{start}$ 的转移安全性证书 $B_{k}(x, q)$，则可以保证 $\mathfrak{L}(\mathfrak{S})$ 中所有\emph{字}的所有运行都不能访问状态 $q_k$。
\end{lemma}

\begin{proof}
  由定义~\ref{transition-safety-certificate-definition}，我们有 $\mathcal{Y}_u = \{(x, q_k) \mid x \in \mathcal{X}\}$ 作为乘积系统 $\mathfrak{S} \times \mathcal{A}$ 中的不安全集。我们验证 $B_k$ 相对于 $\mathcal{Y}_u$ 满足屏障证书条件（\ref{bsinn0}, \ref{bsun0}, \ref{bstnn0}）：
  \begin{itemize}
    \item 条件（\ref{btsc1}）陈述了对所有 $(x_0, q_0) \in \mathcal{Y}_0$，有 $B_k(x_0, q_0) \geq 0$，这对应条件（\ref{bsinn0}）。
    \item 条件（\ref{btsc2}）陈述了对所有 $(x, q_k) \in \mathcal{Y}_u$，有 $B_k(x, q_k) < 0$，这对应条件（\ref{bsun0}）。
    \item 条件（\ref{btsc3}）陈述了对所有 $(x', q') \in F(x, q)$，有 $B_k(x, q) \geq 0 \Rightarrow B_k(x', q') \geq 0$，这对应条件（\ref{bstnn0}）。
  \end{itemize}

  由定理~（\ref{corollary_tbsc}），乘积系统 $\mathfrak{S} \times \mathcal{A}$ 的任意轨迹都不能访问 $\mathcal{Y}_u$。

  现在我们建立与自动机运行之间的联系。对于任意字 $w = \mathfrak{L}(\tau) \in \mathfrak{L}(\mathfrak{S})$，其中 $\tau = \langle x_0, x_1, \dots \rangle$ 是 $\mathfrak{S}$ 的轨迹，$w$ 在 $\mathcal{A}$ 上的对应运行 $r$ 是序列 $r = \langle (r_0, w_0, r_1), (r_1, w_1, r_2), \ldots \rangle$，其中 $r_0 = q_0$。乘积状态序列 $\langle (x_0, r_0), (x_1, r_1), \ldots \rangle$ 构成 $\mathfrak{S} \times \mathcal{A}$ 中的一条轨迹。由于该轨迹不能访问 $\mathcal{Y}_u = \{(x, q_k) \mid x \in \mathcal{X}\}$，我们有对所有 $i$ 成立 $r_i \neq q_k$。因此，$\mathfrak{L}(\mathfrak{S})$ 中所有字的运行都不能访问状态 $q_k$。
\end{proof}

\begin{lemma}
  给定系统 $\mathfrak{S}=(\mathcal{X}, \mathcal{X}_0, f)$、NBA $\mathcal{A} = (\Sigma, Q, q_0, \delta, Acc)$ 和 $(k, l) \in Acc^{pair}$。若能够找到转移持久性证书 $B_{k, l}$，则可以保证 $\mathfrak{L}(\mathfrak{S})$ 中字的所有运行都使从 $q_k$ 到 $q_l$ 的接受转移只发生有限多次。
\end{lemma}

\begin{proof}
我们定义 $\mathcal{X}_{k, l}=\{x \mid (q_k, \mathfrak{L}(x), q_l) \in Acc^{k,l}\}$ 为从 $q_k$ 到 $q_l$ 的接受转移能够发生的状态集合。

考虑任意轨迹 $\tau = \langle x_0, x_1, \dots \rangle \in \mathfrak{T}(\mathfrak{S})$ 及其对应字 $w = \mathfrak{L}(\tau)$。令 $r$ 为 $w$ 在 $\mathcal{A}$ 上从 $q_0$ 出发的任意运行。运行 $r$ 中在步 $i$ 从 $q_k$ 到 $q_l$ 的接受转移需要两个条件：（a）运行在步 $i$ 处于状态 $q_k$，即 $r_i = q_k$；（b）标签使能该转移，即 $(q_k, \mathfrak{L}(x_i), q_l) \in Acc^{k,l}$。由于条件（b）是必要的，只要 $x_i \notin \mathcal{X}_{k,l}$，无论 $r$ 占据哪个自动机状态，从 $q_k$ 到 $q_l$ 的接受转移都不能在步 $i$ 发生。因此，运行 $r$ 中从 $q_k$ 到 $q_l$ 的接受转移次数由轨迹访问 $\mathcal{X}_{k,l}$ 的次数从上方界定，并且这一界对 $w$ 在 $\mathcal{A}$ 上的任意运行 $r$ 都成立。

由定义~\ref{transition-persistence-certificate-definition}，我们有：
\begin{itemize}
  \item 由条件（\ref{btpc1}），$B_{k,l}(x_0) \geq 0$。
  \item 由条件（\ref{btpc2}），对所有 $i$ 有 $B_{k,l}(x_i) \geq 0$，因为 $B_{k,l}$ 沿轨迹保持非负。
  \item 由条件（\ref{btpc3}），$B_{k,l}(x_i) \geq B_{k,l}(x_{i+1}) + I_{k,l}(x_i)\epsilon$，其中若 $x_i \in \mathcal{X}_{k,l}$ 则 $I_{k,l}(x_i) = 1$，否则为 $0$。
\end{itemize}

假设轨迹在索引 $i_1, i_2, \dots, i_m$ 处访问 $\mathcal{X}_{k,l}$，其中 $m$ 是此类访问的次数。每次访问时，$B_{k,l}$ 至少下降 $\epsilon$：
\begin{align}
B_{k,l}(x_0) \geq B_{k,l}(x_{i_1}) + \epsilon \geq B_{k,l}(x_{i_2}) + 2\epsilon \geq \cdots \geq B_{k,l}(x_{i_m}) + m\epsilon.
\end{align}

由于由条件（\ref{btpc2}）有 $B_{k,l}(x_{i_m}) \geq 0$，于是 $B_{k,l}(x_0) \geq m\epsilon$，从而 $m \leq \frac{B_{k,l}(x_0)}{\epsilon}$。由于任意运行中从 $q_k$ 到 $q_l$ 的接受转移次数都由 $m$ 从上方界定，因此它们只发生有限多次。由于该界仅依赖于系统轨迹，而不依赖于运行的选择，所以该结论对 $\mathfrak{L}(\mathfrak{S})$ 中所有字的所有运行均成立。
\end{proof}


\section{实验}
\label{appendix-experiment}
\subsection{二维 Kuramoto 振子系统的模板}
二维 Kuramoto 振子系统的转移安全性证书模板给定为：
\begin{align*}
B_0(x_1, x_2, q) = &\, c_0 + c_1 \cdot x_1 + c_2 \cdot x_2 + c_3 \cdot q \\
                   &+ c_4 \cdot I_{X_u}(x_1, x_2) + c_5 \cdot x_1 q + c_6 \cdot x_2 q
\end{align*}

\subsection{双房间温度系统的模板}
对于\textit{转移安全性证书}：
\begin{align*}
B_1(x_1, x_2, q) = &\, c_0 + c_1 \cdot I_{X_0}(x_1, x_2) + c_2 \cdot I_{X_{VF}}(x_1, x_2) \\
                   &+ c_3 \cdot x_1 I_{X_0}(x_1, x_2) + c_4 \cdot x_2 I_{X_0}(x_1, x_2) \\
                   &+ c_5 \cdot x_1 I_{X_{VF}}(x_1, x_2) + c_6 \cdot x_2 I_{X_{VF}}(x_1, x_2) \\
                   &+ c_7 \cdot \max(x_1, x_2) + c_8 \cdot x_1^2 + c_9 \cdot x_2^2 + c_{10} \cdot q + c_{11} \cdot q^2.
\end{align*}
\textit{转移持久性证书}的模板如下：
\begin{align*}
B_{1,1}(x_1, x_2) = & \, c_0 + c_1 \cdot I_{X_0}(x_1, x_2) + c_2 \cdot I_{X_{VF}}(x_1, x_2) \\
                    & + c_3 \cdot x_1 I_{X_0}(x_1, x_2) + c_4 \cdot x_2 I_{X_0}(x_1, x_2) \\
                    & + c_5 \cdot x_1 I_{X_{VF}}(x_1, x_2) + c_6 \cdot x_2 I_{X_{VF}}(x_1, x_2) \\ 
                    & + c_7 \cdot \max(x_1, x_2) + c_8 \cdot x_1^2 + c_9 \cdot x_2^2.
\end{align*}

\subsection{神经网络证书合成结果}
\label{app:neural_cert_results}

本节给出了三个案例研究中合成得到的神经网络证书。

	\subsubsection{Kuramoto 振子：一维系统}
	\label{app:ko1d_neural}
	\textbf{有效转移安全性证书网络参数：}
	\begin{itemize}
	    \item \textbf{网络架构：} $2 \rightarrow 10 \rightarrow 1$（ReLU 激活网络）
	    \item \textbf{第 1 层权重} $W^{(1)} \in \mathbb{R}^{10 \times 2}$：
	    \[
	    W^{(1)} = \begin{bmatrix}
	     0.35423541 &  0.03017079 \\
	     0.40347892 &  0.15485710 \\
	    -0.33290595 &  0.20670444 \\
	     0.05953831 & -0.57081002 \\
	    -0.77764583 & -0.30270880 \\
	     0.28167149 & -0.56571352 \\
	    -0.00433345 &  0.27621469 \\
	    -0.11347809 & -0.37141919 \\
	     0.30172166 & -0.14517429 \\
	     0.00136681 &  0.32282710
	    \end{bmatrix}
	    \]
	    
	    \item \textbf{第 1 层偏置} $b^{(1)} \in \mathbb{R}^{10}$：
	    \[
	    b^{(1)} = \begin{bmatrix}
	    -0.83490396 &  0.57956743 & -0.47568363 & -0.27932465 & -0.28890583 \\
	     0.16922584 & -0.66955644 & -0.23689666 &  0.49750403 &  0.55335736
	    \end{bmatrix}^\top
	    \]
	    
	    \item \textbf{第 2 层权重} $W^{(2)} \in \mathbb{R}^{1 \times 10}$：
	    \[
	    W^{(2)} = \begin{bmatrix}
	    -1.36676764 &  0.13470453 & -0.01585477 & -0.04601728 &  0.22054167 \\
	    -0.15794340 & -0.09915265 &  0.00324699 & -0.02935893 &  0.18765648
	    \end{bmatrix}
	    \]
	    
	    \item \textbf{第 2 层偏置} $b^{(2)} \in \mathbb{R}$：
	    \[
	    b^{(2)} = -0.25243226
	    \]
	    
	    \item \textbf{解析表达式：}
	    \[
	    B(x, q) = W^{(2)} \cdot \text{ReLU}(W^{(1)} \cdot [x, q]^\top + b^{(1)}) + b^{(2)}
	    \]
	\end{itemize}

	\subsubsection{Kuramoto 振子：二维系统}
	\label{app:ko2d_neural}
	
	\textbf{有效转移安全性证书网络参数：}
	
	\begin{itemize}
	    \item \textbf{网络架构：} $3 \rightarrow 15 \rightarrow 1$（ReLU 激活网络）
	    \item \textbf{第 1 层权重} $W^{(1)} \in \mathbb{R}^{15 \times 3}$：
	    {\footnotesize
	    \[
	    W^{(1)} = \begin{bmatrix}
	    -0.01797801 &  0.64069444 & -0.03856681 \\
	    -0.57299656 & -0.14924462 & -0.04329325 \\
	    -0.56604266 & -0.25653100 &  0.10918397 \\
	    -0.12088224 &  0.62192768 & -0.95162183 \\
	    -0.54926884 & -0.45272782 &  0.26691911 \\
	    -0.35724673 &  0.10176215 & -0.26044405 \\
	    -0.62156856 &  0.45307916 & -0.63630158 \\
	    -0.50465399 & -0.00161430 &  0.04746520 \\
	    -0.21175091 &  0.12297605 & -0.52375120 \\
	     0.60632193 &  0.14016706 & -2.34126782 \\
	    -0.07868959 & -0.15929390 & -0.10671024 \\
	     0.68689883 & -0.58823287 & -0.65129077 \\
	    -0.75197023 &  0.13900414 & -0.01829747 \\
	    -0.55132180 &  0.17512728 & -0.04693245 \\
	    -0.18292713 &  0.10734314 & -0.20697621
	    \end{bmatrix}
	    \]}
	    
	    \item \textbf{第 1 层偏置} $b^{(1)} \in \mathbb{R}^{15}$：
	    {\scriptsize
	    \[
	    b^{(1)} = \begin{bmatrix}
	    -0.23332441 & -0.25264311 & -0.20304658 &  0.00498547 & -0.41519192 \\
	    -0.18106091 & -0.32431474 &  0.83221281 & -0.54517263 &  0.07074507 \\
	    -0.18164203 &  0.44940028 &  0.60892761 & -0.63835907 &  0.36881346
	    \end{bmatrix}^\top
	    \]}
	    
	    \item \textbf{第 2 层权重} $W^{(2)} \in \mathbb{R}^{1 \times 15}$：
	    {\footnotesize
	    \[
	    W^{(2)} = \begin{bmatrix}
	    -0.22106327 &  0.12717371 & -0.04997402 & -0.50930381 & -0.18251525 \\
	     0.08938225 &  0.48234484 & -0.43327656 & -0.12058055 &  0.42470711 \\
	     0.11804953 & -0.65379494 &  0.25955433 &  0.05265374 &  0.01196656
	    \end{bmatrix}
	    \]}
	    
	    \item \textbf{第 2 层偏置} $b^{(2)} \in \mathbb{R}$：
	    \[
	    b^{(2)} = 0.36412391
	    \]
	    
	    \item \textbf{解析表达式：}
	    \[
	    B(x_1, x_2, q) = W^{(2)} \cdot \text{ReLU}(W^{(1)} \cdot [x_1, x_2, q]^\top + b^{(1)}) + b^{(2)}
	    \]
	\end{itemize}

\subsubsection{双房间温度系统}
\label{app:two_room_neural}
\textbf{有效转移持久性证书网络参数：}

\begin{itemize}
    \item \textbf{网络架构：} $2 \rightarrow 3 \rightarrow 1$
    \item \textbf{下降常数 $\epsilon$：} 0.01（固定）
    \item \textbf{第 1 层权重} $W^{(1)} \in \mathbb{R}^{3 \times 2}$：
    \[
    W^{(1)} = \begin{bmatrix}
    -0.12892260 &  0.07935464 \\
     0.39971098 & -0.35251060 \\
    -0.61775041 &  0.61788529
    \end{bmatrix}
    \]
    
    \item \textbf{第 1 层偏置} $b^{(1)} \in \mathbb{R}^{3}$：
    \[
    b^{(1)} = \begin{bmatrix}
     1.44339669 & -1.41302800 & -0.01615771
    \end{bmatrix}^\top
    \]
\end{itemize}

\subsection{比较方法设置}

\subsubsection{状态三元组方法}
状态三元组基线遵循~\cite{wongpiromsarn2015automata} 的接受状态 NBA 路径阻断方案。对于每个接受状态，枚举从初始自动机状态到该接受状态的所有简单路径。对于三元组 $(q_i,q_j,q_k)$，令 $Y_0$ 为 $q_i\rightarrow q_j$ 上标签的并集，$Y_1$ 为 $q_j\rightarrow q_k$ 上标签的并集，而当存在自环 $q_j\rightarrow q_j$ 时，$Y_m$ 为该自环上的标签。我们合成满足以下条件的屏障 $B$：
\[
B(x)\le 0 \text{ 在 } Y_0\text{ 上},\qquad
B(x)\ge 10^{-2} \text{ 在 } Y_1\text{ 上},
\]
以及不变性条件
\[
x\in (Y_0\cup Y_m\cup Y_1)\setminus Y_1,\; f(x)\in Y_0\cup Y_m\cup Y_1,\; B(x)\le 0
\Rightarrow B(f(x))\le 0.
\]
候选系数由 Z3 合成，并针对 Kuramoto 系统由 dReal 形式化检查，针对双房间系统由 Z3 形式化检查。

对于一维 Kuramoto 基准，状态三元组模板为
\[
B(x)=c_0+c_1x+c_2x^2+c_3x^3+c_4x^4.
\]
合成得到的证书为
\[
B(x)=-4.5x^2+2x^3,
\]
该证书阻断了接受路径 $0\rightarrow1\rightarrow2$，其中 $Y_0=\{i\}$、$Y_1=\{u\}\cup\{i,u\}$、$Y_m=\emptyset\cup\{i\}$。直接接受边 $0\rightarrow2$ 因其标签区域 $\{i,u\}$ 为空而被处理。

对于二维 Kuramoto 基准，测试的状态三元组模板为
\[
B(x_1,x_2)=c_0+c_1x_1+c_2x_2+c_3x_1x_2+c_4x_1^2+c_5x_2^2.
\]
直接接受边 $0\rightarrow2$ 同样由于空标签区域 $\{i,u\}$ 而被处理。路径 $0\rightarrow1\rightarrow2$ 被如下合成证书阻断：
\[
B(x_1,x_2)=-x_1-x_2-3.125x_1x_2+2x_1^2+2x_2^2,
\]
其中 $Y_0=\{i\}$、$Y_1=\{u\}\cup\{i,u\}$、$Y_m=\emptyset\cup\{i\}$。该运行对依赖动力学的不变性检查使用了 dReal 精度 $10^{-3}$。

对于双房间基准，测试的状态三元组模板为
\[
B(x_1,x_2)=c_0+c_1x_1+c_2x_2+c_3x_1x_2+c_4x_1^2+c_5x_2^2+c_6\max(x_1,x_2).
\]
接受路径 $0\rightarrow2\rightarrow1$ 未在超时前被处理，其中 $Y_0=\{a\}\cup\{a,b\}$、$Y_1=\{b\}\cup\{a,b\}$、$Y_m=\emptyset\cup\{a\}$。对于这一持久性风格的基准，这是预期现象：由于 $\mathcal{X}_0\subseteq \mathcal{X}_{VF}$，来自初始区域的轨迹可能同时满足 $a$ 和 $b$，因此基于状态的 NBA 的接受状态可达，而状态三元组的路径阻断义务过于保守。

\subsubsection{闭包证书方法}
合成过程使用 Z3 寻找候选证书，并使用 dReal 寻找反例，其中 dReal 精度为 0.0001。

\textbf{一维 Kuramoto 振子系统：}
NBA 定义为 $(\Sigma, Q, q_0, \delta, \text{Acc})$，用于表示 $F a$。
\begin{align*}
    \text{Acc} &= \{0\} \\
    Q &= \{1, 0\} \\
    q_0 &= 1 \\
    \Sigma &= \{\emptyset, \{a\}\} \\
    \delta &= \{(1, \emptyset, 1), (1, \{a\}, 0), (0, \emptyset, 0), (0, \{a\}, 0)\} \\
	a &\text{ 定义为 } x \in X_{\text{u}}
\end{align*}
证书模板为：
\begin{align*}
    T_{0,0} &= c_0 + c_1 x + c_2 y \\
    T_{0,1} &= c_3 + c_4 x + c_5 y \\
    T_{1,0} &= c_6 + c_7 x + c_8 y \\
    T_{1,1} &= c_9 + c_{10} x + c_{11} y
\end{align*}
最终合成得到的系数见表~\ref{tab:coeff_1d}。

\begin{table}[h]
\centering
\caption{合成闭包证书的系数（一维系统）}
\label{tab:coeff_1d}
\begin{tabular}{lccc}
\hline
\textbf{模板} & $c_0$ & $c_1$ & $c_2$ \\
\hline
$T_{0,0}(x, y)$ & 0.000000 & 0.000000 & 0.000000 \\
$T_{0,1}(x, y)$ & -0.500000 & 0.000000 & 0.000000 \\
$T_{1,0}(x, y)$ & -4.310550 & 1.764117 & 0.000000 \\
$T_{1,1}(x, y)$ & 0.000000 & 0.000000 & 0.000000 \\
\hline
\end{tabular}
\end{table}

\textbf{二维 Kuramoto 振子系统：}
在该模板下，CEGIS 循环超时（超过一小时），这与~\cite{murali2024closure} 中报告的结果一致，其中在原设置下二维基准的合成大约需要 1 小时 50 分钟。
证书模板对 $i, j \in \{0, 1\}$ 定义为：
\begin{align*}
T_{i,j}((x_1, x_2), (y_1, y_2)) = &\, c_{0, i, j} + c_{1, i, j} y_1 I_{X_0}(x_1, x_2) + c_{2, i, j} y_2 I_{X_0}(x_1, x_2) \\
                                  &+ c_{3, i, j} y_1 I_{X_u}(x_1, x_2) + c_{4, i, j} y_2 I_{X_u}(x_1, x_2) \\
                                  &+ c_{5,i,j} y_1 + c_{6,i,j} y_2
\end{align*}

\textbf{双房间温度系统：}
在该模板下，CEGIS 循环超时（超过一小时），这与~\cite{murali2024closure} 中报告的结果一致，其中在原设置下该基准的合成需要超过 1.5 小时。
用于 $a_0 \land F a_1$ 的 NBA 为 $(\Sigma, Q, q_0, \delta, \text{Acc})$。

\begin{align*}
    Q &= \{0,1,2,3\} \\
    \Sigma &= \{\emptyset, \{a_0\}, \{a_1\}, \{a_0, a_1\}\} \\
    q_0 &= 0 \\
    \text{Acc} &= \{2\} \\
    \delta &= \{(q_0, \{a_0\}, q_1), (q_0, \{a_1\}, q_1), (q_0, \{a_0, a_1\}, q_2), (q_0, \emptyset, q_3), \\
    &\quad\;(q_0, \{a_1\}, q_3), (q_1, \emptyset, q_1), (q_1, \{a_0\}, q_1), (q_1, \{a_1\}, q_2), \\
    &\quad\;(q_1, \{a_0, a_1\}, q_2), (q_2, \emptyset, q_2), (q_2, \{a_0\}, q_2), (q_2, \{a_1\}, q_2), \\
    &\quad\;(q_2, \{a_0, a_1\}, q_2), (q_3, \emptyset, q_3), (q_3,\{a_0\},q_3), (q_3, \{a_1\}, q_3), \\
    &\quad\;(q_3, \{a_0, a_1\}, q_3)\} \\
	a_0 &\text{ 定义为 } x \in X_{0} \\
	a_1 &\text{ 定义为 } x \in X_{VF} \\
\end{align*}
该系统的证书模板对 $i, j \in \{0, 1, 2, 3\}$ 定义为：
\begin{align*}
T_{i,j}((x_1, x_2), (y_1, y_2)) = &\, c_{0, i, j} + c_{1, i, j}x_1 + c_{2, i, j}x_2 + c_{3, i, j}y_1 + c_{4, i, j}y_2 \\
                                  &+ c_{5, i, j}\max(x_1, x_2) + c_{6, i, j}\max(y_1, y_2) \\
                                  &+ c_{7, i, j}x_1^2 + c_{8, i, j}x_2^2 + c_{9, i, j}y_1^2 + c_{10, i, j}y_2^2
\end{align*}

\subsubsection{神经闭包证书方法}
神经闭包证书的配置总结见表~\ref{tab:ncc_config}。所有系统均使用采样密度 $\xi$ 为 0.01，安全裕度 $\eta$ 为 0.01。Lipschitz 常数使用基于谱范数的算法估计。

\begin{table}[h]
\centering
\caption{神经闭包证书的神经网络架构与参数}
\label{tab:ncc_config}
\begin{tabular}{lccc}
\hline
\textbf{系统} & \textbf{网络架构} & \textbf{激活} & \textbf{Lipschitz 估计} \\
\hline
一维 Kuramoto 振子 & (4-80-1) & ReLU & 谱范数 \\
二维 Kuramoto 振子 & (6-80-1) & ReLU & 谱范数 \\
双房间温度 & (6-80-1) & ReLU & 谱范数 \\
\hline
\end{tabular}
\end{table}
